% IEEE-style write-up for RAG-based threat modeling and attack vector generation

\subsection{Agentic RAG for Automated Threat Intelligence}

The ARES framework employs an agentic Retrieval-Augmented Generation (RAG) system to automate the discovery and classification of cybersecurity vulnerabilities across the EVCS-integrated transmission-distribution infrastructure. The RAG agent is built on the Gemini 2.5 Flash foundation model, augmented with a domain-specific knowledge base comprising 424 documents spanning multiple vulnerability intelligence sources. The knowledge base integrates structured threat taxonomies (MITRE ATT\&CK for ICS, STRIDE patterns), real-world vulnerability disclosures (NVD CVE database, ICS-CERT advisories), protocol-specific security analyses (OCPP, DNP3, Modbus, TCP/IP), empirical attack datasets (CICEVSE2024), and peer-reviewed research literature on EVCS and power grid security. This multi-source integration enables the RAG agent to ground its threat assessments in both theoretical frameworks and empirically observed attack patterns.

The knowledge base construction pipeline processes each source through specialized collectors that extract vulnerability metadata, map threats to MITRE ATT\&CK techniques and STRIDE categories, identify affected systems and protocols, and generate semantic embeddings using the all-MiniLM-L6-v2 sentence transformer model. The resulting 384-dimensional embeddings are indexed in a ChromaDB vector database, enabling efficient semantic retrieval of contextually relevant threat intelligence for arbitrary queries. When presented with a vulnerability discovery query, the RAG agent retrieves the top-k most semantically similar documents from the knowledge base, constructs a context-augmented prompt that explicitly instructs the foundation model to extract CVE identifiers, MITRE technique IDs, and STRIDE categories from the retrieved context, and generates a structured threat assessment that maps high-level threat categories to concrete attack vectors with protocol-specific implementation details.

To evaluate the quality and reliability of RAG-generated threat intelligence, we developed a confidence scoring mechanism that assesses six dimensions of response quality: CVE citation accuracy (25 points), MITRE technique coverage (25 points), RL action specificity (20 points), protocol specificity (15 points), context document utilization (10 points), and structured format adherence (5 points). This scoring function enables automated selection of high-confidence vulnerability mappings for subsequent RL-based attack simulation. Comparative evaluation against a non-RAG baseline demonstrated that knowledge base augmentation improves threat intelligence quality by an average of 32.9 confidence points across 15 structured discovery queries, with RAG outperforming the baseline in 14 of 15 queries. The RAG agent achieved confidence scores ranging from 50.0 to 92.0 across different STRIDE categories, with Denial of Service and Spoofing vulnerabilities exhibiting the strongest empirical support due to extensive CVE documentation in the knowledge base.

\subsection{STRIDE-MITRE Integration for Structured Attack Modeling}

STRIDE and MITRE ATT\&CK for ICS provide complementary perspectives on cyber threat analysis for critical infrastructure. STRIDE, developed by Microsoft, offers a systematic threat categorization framework that maps six threat types—Spoofing, Tampering, Repudiation, Information Disclosure, Denial of Service, and Elevation of Privilege—to fundamental security properties including authentication, integrity, non-repudiation, confidentiality, availability, and authorization. This architecture-centric methodology excels at design-phase threat enumeration, enabling systematic identification of attack surfaces through data flow diagram analysis. MITRE ATT\&CK for ICS, in contrast, provides an empirically grounded taxonomy of adversary tactics, techniques, and procedures derived from real-world cyber incidents affecting operational technology environments, including notable attacks such as Stuxnet, BlackEnergy, and Industroyer. The framework catalogs 81 techniques across 12 tactical categories, offering behavioral models of adversary actions that inform detection and defense strategies.

The ARES framework explicitly integrates these complementary frameworks through agentic RAG-based mapping. The RAG agent applies STRIDE threat modeling to the EVCS-integrated transmission-distribution architecture, systematically enumerating component-level threats across nine communication links spanning the cyber-physical interface from electric vehicles to automatic generation control. For each identified STRIDE threat, the agent queries the knowledge base to retrieve contextually relevant MITRE ATT\&CK techniques, producing a structured mapping of the form (system element, STRIDE category) → (MITRE ATT\&CK tactic, technique IDs). This mapping bridges high-level architectural threat modeling with concrete adversary behavior patterns, enabling the translation of abstract threat categories into actionable attack scenarios for RL-based exploration.

The integration process leverages the RAG agent's ability to synthesize information across heterogeneous knowledge sources. When analyzing a specific communication link—for example, the OCPP-based interface between electric vehicles and charging stations—the agent retrieves STRIDE threat patterns specific to authentication and message integrity, ICS-CERT advisories documenting real-world OCPP vulnerabilities with assigned CVE identifiers, protocol security analyses detailing cryptographic weaknesses and implementation flaws, and MITRE techniques such as T0855 (Unauthorized Command Message) and T0866 (Manipulation of View) that describe adversary actions exploiting these vulnerabilities. The resulting threat assessment connects architectural vulnerabilities to empirically observed attack techniques, grounding the threat model in both design-phase analysis and operational intelligence.

\subsection{Data Flow-Based Threat Assessment}

The threat modeling process begins with the construction of a comprehensive data flow diagram representing the EVCS-integrated transmission-distribution architecture. The architecture comprises two interconnected layers: a physical layer spanning electric vehicles, charging stations, distribution grid infrastructure, transmission grid components, and generation assets; and a cyber layer encompassing charging management systems (CMS), distribution system management (DSM), energy management systems (EMS), and automatic generation control (AGC). Nine primary data flows traverse this architecture, each characterized by specific protocols, trust boundaries, and operational semantics.

Data Flow 1 carries charging information including state of charge, power demand, and optimal reference values for voltage, current, and power between electric vehicles and charging stations using the Open Charge Point Protocol (OCPP). This bidirectional communication enables dynamic charging control but introduces authentication and message integrity vulnerabilities. Data Flows 2 through 6 represent the hierarchical aggregation of load measurements and control signals from local charging management through distribution system management to energy management systems, primarily using DNP3 and TCP/IP protocols. Data Flows 7 through 9 complete the control loop, carrying frequency measurements and optimal reference setpoints from the transmission grid through EMS to AGC, enabling grid-wide power and frequency stabilization.

The RAG agent systematically analyzes each data flow to identify STRIDE threats and map them to MITRE ATT\&CK techniques. For Data Flow 1, the agent identifies Spoofing threats arising from weak OCPP authentication mechanisms, mapping them to MITRE techniques T0855 (Unauthorized Command Message) and T0862 (Rogue Device), with supporting evidence from CVE-2024-23971 and ICS-CERT advisory ICSA-22-256-01. For Data Flows 3 and 4, which carry unencrypted DNP3 load measurements between CMS and distribution grid components, the agent identifies Information Disclosure threats mapped to T0842 (Network Traffic Monitoring) and Tampering threats mapped to T0836 (Modify Parameter), supported by protocol vulnerability analyses in the knowledge base. This systematic, flow-by-flow analysis ensures comprehensive coverage of the attack surface while maintaining traceability between architectural elements, threat categories, and adversary techniques.

\subsection{STRIDE-MITRE-RL Action Mapping}

The culmination of the threat modeling process is the generation of a structured mapping between STRIDE categories, MITRE ATT\&CK techniques, and concrete RL attack actions suitable for simulation in the federated PINN digital twin. This mapping serves as the interface between high-level threat intelligence and low-level attack implementation, translating abstract vulnerability categories into executable attack strategies that manipulate specific state variables in the cyber-physical simulation environment. Table~\ref{tab:stride_mitre_rl_detailed} presents the complete mapping, derived from the confidence-ranked outputs of the agentic RAG evaluation.

Spoofing vulnerabilities, identified with the highest confidence score of 92.0, map to Communication Spoofing attacks targeting Data Flow 1 (EV ↔ EVCS via OCPP). The RAG agent identified specific attack vectors including spoofing of OCPP Authorize.conf messages to permit unauthorized charging, StopTransaction.req messages to prematurely terminate sessions, and MeterValues.req messages to submit falsified billing data. The corresponding MITRE techniques—T0817 (Drive-by Compromise), T0819 (Exploit Public-Facing Application), and T0867 (Lateral Tool Transfer)—describe the adversary actions required to gain initial access and impersonate legitimate entities. In the RL simulation environment, this translates to an attack action that injects spoofed OCPP messages into the communication channel, causing the charging management system to accept commands from unauthorized sources.

Tampering vulnerabilities map to Data Injection attacks targeting Data Flows 3 and 4 (CMS ↔ Distribution Grid via DNP3), where the RL agent injects fabricated analog input values including false grid frequency deviations and manipulated demand factor measurements. The RAG agent identified this attack vector with confidence scores of 57.0 to 73.0, supported by MITRE techniques T0806 (Brute Force I/O), T0836 (Modify Parameter), and T0871 (Execution through API). The attack exploits the integrity vulnerabilities in DNP3 protocol implementations to corrupt load measurement data feeding the CMS PINN optimizer, causing it to compute incorrect charging references based on falsified grid conditions. This directly violates system integrity by inserting unauthorized modifications into the control data stream.

Repudiation vulnerabilities map to Protocol Manipulation attacks targeting Data Flow 5 (DSM ↔ EMS via DNP3) and OCPP-based flows. The RAG agent identified these vulnerabilities with confidence scores of 55.0 to 73.0, noting the absence of Secure Authentication v5 in many DNP3 implementations and the lack of immutable audit trails in OCPP systems. The corresponding MITRE technique T0872 (Indicator Removal on Host) describes adversary actions to erase evidence of malicious activity. In the simulation, the RL agent manipulates protocol-level features including forged timestamps, suppressed acknowledgments, and unsolicited responses to inject false load forecasts while maintaining plausible deniability, enabling the attacker to deny responsibility for malicious actions.

Information Disclosure vulnerabilities map to Voltage Manipulation attacks through a two-phase strategy. In the reconnaissance phase, the RL agent exploits unencrypted DNP3 communication on Data Flows 3 and 4 to passively intercept load and voltage measurement data, learning real-time grid voltage levels, capacity limits, and operational patterns. This corresponds to MITRE techniques T0840 (Network Connection Enumeration), T0802 (Automated Collection), and T0842 (Network Traffic Monitoring). In the exploitation phase, the agent uses the acquired knowledge to craft precise voltage manipulation attacks calibrated to remain below alarm thresholds, evading simple detection mechanisms. The RAG agent identified this attack vector with confidence scores of 50.0 to 78.0, demonstrating that information disclosure enables more sophisticated follow-on attacks.

Denial of Service vulnerabilities, identified with the highest aggregate confidence scores (81.0 to 92.0), map to Power Disruption attacks targeting Data Flows 4 and 9. The RAG agent identified extensive CVE evidence for DoS attacks on both DNP3 and TCP/IP protocols, including TCP SYN flooding, UDP flooding, DNP3 unsolicited response flooding, and malformed packet injection causing buffer overflows or daemon crashes. The corresponding MITRE techniques—T0800 (Activate Firmware Update Mode), T0814 (Denial of Service), and T0816 (Device Restart/Shutdown)—describe adversary actions to disrupt system availability. In the simulation, disruption of Data Flow 9 prevents AGC from receiving optimal reference setpoints, impairing grid-wide frequency stabilization, while disruption of Data Flow 4 blinds distribution system management from real-time load visibility.

Elevation of Privilege vulnerabilities map to Current Injection attacks targeting Data Flows 1 and 2 (EV ↔ EVCS ↔ CMS). The RAG agent identified privilege escalation vulnerabilities including hard-coded credentials (CVE-2021-22707), firmware exploitation, and insecure API endpoints that enable attackers to gain administrative access and modify charging current and voltage limits. The corresponding MITRE techniques—T0890 (Exploitation for Privilege Escalation), T0891 (Hardcoded Credentials), and T0839 (Module Firmware)—describe the adversary actions required to escalate from network user to administrator privileges. In the simulation, the RL agent exploits weak authentication on Data Flow 2 to escalate privileges, then overrides current constraints on Data Flow 1 to drive unsafe overcurrent conditions that violate equipment ratings and protection relay settings.

This structured STRIDE-MITRE-RL mapping provides a systematic translation from high-level threat modeling through empirically grounded adversary techniques to executable attack actions in the simulation environment. Each mapping is supported by specific CVE identifiers, MITRE technique IDs, knowledge base references, and confidence scores derived from the agentic RAG evaluation, ensuring traceability and empirical grounding. The mapping collectively spans three protocol families (OCPP, DNP3, TCP/IP), six STRIDE categories, and six primary communication links in the EVCS-to-grid architecture, providing comprehensive coverage of the cyber-physical attack surface. The confidence-based ranking mechanism ensures that the selected attack actions represent the most well-documented and empirically supported vulnerabilities, maximizing the realism and relevance of the RL-based attack exploration in the federated PINN digital twin.
